\section{Decay--Delay Analysis}
\label{sec:theory}

This section derives a closed-form relationship between a channel's decay
rate and (a) how long the model takes to react to an anomaly onset and
(b) the shortest anomaly it can reliably detect, then validates it
empirically (Section~\ref{sec:experiments}, via
\texttt{scripts/analyze\_theory.py}). The analysis treats one channel of
the recurrence (Eq.~\ref{eq:recurrence}) in isolation with its
\emph{effective} decay $a$ held constant across the transition being
analyzed; Section~\ref{sec:theory-gating} discusses how the learned gate
$g_t$ relaxes this assumption in practice.

\subsection{Settling-time bound}
\label{sec:theory-settling}
Consider a single channel driven by an input that is constant at
$u_{\text{norm}}$ for $t < t_0$ and step-changes to $u_{\text{anom}}$ at
$t_0$ (an idealized anomaly onset). Under the recurrence
$s_t = a\, s_{t-1} + u_t$, the state converges toward the steady state
$s^\ast = u/(1-a)$ for constant $u$; writing $s_{\text{norm}}^\ast =
u_{\text{norm}}/(1-a)$ and $s_{\text{anom}}^\ast = u_{\text{anom}}/(1-a)$,
the solution for $t \ge t_0$ is
\begin{equation}
  s_t - s_{\text{anom}}^\ast = a^{\,t-t_0}\big(s_{t_0} - s_{\text{anom}}^\ast\big).
  \label{eq:transient}
\end{equation}

\begin{proposition}[Settling delay]
\label{prop:settling}
Define the detection delay $\delta(\varepsilon)$ as the smallest
$t - t_0 \ge 0$ such that the state has moved to within a fraction
$\varepsilon \in (0,1)$ of the new steady state, i.e.
$|s_t - s_{\text{anom}}^\ast| \le \varepsilon\, |s_{t_0} -
s_{\text{anom}}^\ast|$. Then
\begin{equation}
  \delta(\varepsilon) = \left\lceil \frac{\ln \varepsilon}{\ln a} \right\rceil.
  \label{eq:delay-bound}
\end{equation}
\end{proposition}
\begin{proof}
From Eq.~\ref{eq:transient}, the condition is $a^{\,\delta} \le
\varepsilon$. Since $a \in (0,1)$, $\ln a < 0$; taking logarithms and
dividing by $\ln a$ flips the inequality, giving $\delta \ge \ln
\varepsilon / \ln a$ (positive, since $\ln\varepsilon < 0$ too). The
smallest integer $\delta$ satisfying this is
Eq.~\ref{eq:delay-bound}.
\end{proof}

Eq.~\ref{eq:delay-bound} is monotonically increasing in $a$: channels
with slower decay (larger $a$, longer memory) take strictly longer to
react to a step change. This is exactly the quantity
\texttt{scripts/analyze\_theory.py} computes from each trained layer's
mean learned base decay $\bar{a}$ (Eq.~\ref{eq:base-decay}) and compares
against the empirically measured delay between an anomaly's labeled onset
and the first score-threshold crossing.

\subsection{Minimum-detectable event duration}
\label{sec:theory-min-duration}
Proposition~\ref{prop:settling} has a direct converse: an anomalous
segment of duration $D$ frames that ends before the state has settled
(i.e. $D < \delta(\varepsilon)$ for the $\varepsilon$ at which the
downstream threshold reliably separates normal/anomalous scores) will
not produce the full-magnitude prediction error the detector was
calibrated on, and is at elevated risk of a missed detection. This gives
a direct, checkable prediction: \emph{shortening} the base decay (smaller
$a_{\min}, a_{\max}$ in Eq.~\ref{eq:base-decay}) should reduce delay and
improve recall on short anomalous segments, at the cost of a noisier
score curve (and likely higher false-positive rate) on normal video,
since the state now also reacts to high-frequency variation that is not
anomalous. Section~\ref{sec:experiments} reports the state-size/decay
sweep ablation that tests this trade-off directly.

\subsection{Role of the event-boundary gate}
\label{sec:theory-gating}
The analysis above assumes $a$ is fixed, but the actual effective decay
$a_t = \bar a \odot g_t$ (Eq.~\ref{eq:gate}) is state- and
input-dependent. Qualitatively, this was intended to let the model behave
like a \emph{small} $a$ (fast reaction, per Eq.~\ref{eq:delay-bound})
exactly when $g_t$ drops in response to a state/input mismatch — i.e.\ at
a genuine event boundary — while behaving like the larger, more stable
base decay $\bar a$ elsewhere: a shorter effective delay than a
fixed-decay recurrence with the same $\bar a$, without paying the
false-positive cost of a globally smaller decay everywhere.

The settling-delay numbers in Section~\ref{sec:experiments} are
consistent with the gate reacting fast (empirical delay far below the
fixed-$\bar a$ bound on both Ped2 and Avenue). However, the ablations in
Tables~\ref{tab:ablation-ucsd-ped2}--\ref{tab:ablation-cuhk-avenue}
complicate the story, and in a way that is itself informative: fast
reaction to a state/input mismatch is not automatically the same as
reacting specifically to \emph{anomalous} mismatches, and whether the
gate's extra parameters (Eq.~\ref{eq:gate}) learn the latter rather than
overfitting to the former appears to depend on how much training data is
available. On Ped2 (16 clips, 120-180 frames each — the smaller
training set), disabling the gate \emph{improves} frame-AUC (76.6\% vs.\
67.9\%), and a smaller state does too — both point to overfitting. On
Avenue (16 clips, but up to 1271 frames each — substantially more total
training frames), the pattern reverses: disabling the gate \emph{hurts}
sharply (60.0\% vs.\ 70.2\%), and a larger state helps slightly rather
than hurting. Read together, this is more consistent with a
capacity/data-size interaction than with the gating mechanism being
unhelpful in principle — but it is a two-dataset hypothesis, not a
demonstrated effect, and a third, larger training set
(ShanghaiTech, in progress) is needed before treating it as settled. We
report the full, dataset-dependent picture here rather than only the
settling-delay result that would have supported the original hypothesis
on its own.
% TODO: extend Proposition 1 to the time-varying-$a_t$ case (closed form
% via a cumulative product, see the parallel-form discussion in
% src/models/ssm.py's docstring) for a tighter bound; and re-run the gate
% ablation on Avenue/ShanghaiTech (more training clips) before drawing a
% dataset-independent conclusion.
